\scanpage{81}
Ainsi
\[
X\unsure{-|\longrightarrow|\longmapsto}\Gamma^*(X)=(\Gamma^i(X))_{i\geq0}
\]
est un foncteur (pas \inserted{\unsure{néct|ncés}} additif !)
\[
P_1\longrightarrow\operatorname{Grad}(P_1)
\]
qui est en fait un foncteur
``transformant sommes en $\otimes$''. Comme
tout objet de $P_1$ est un \inserted{objet} groupe comm.
pour l'opération $\oplus$, il s'ensuit que
dans chaque $\Gamma^*(X)$ il y a une structure
de $\otimes$-algèbre \inserted{graduée} comm., fonctorielle en $X$.
% slanted
\margin{Pour la $=$ raison $\Gamma^*X$ sera \unsure{est} aussi une cogèbre, \deleted{\unsure{donc}} et $=$ une \uline{bigèbre} ass. commut pour les $2$ structures.}

\deleted{En particulier}
On va définir \inserted{ds $P_1$} un hom, fonctoriel en $X$
\[
\Gamma^p(\Gamma^q(X))\longleftarrow\Gamma^{pq}(X)\tag*{($*$)}
\]
ou ce qui revient au $=$, un hom \deleted{\unsure{dans}}
de degré $pq$
\[
X\longrightarrow\Gamma^p(\Gamma^q(X))
\]
On prendra le composé
\diagram{\begin{tikzcd}[ampersand replacement=\&,column sep=large]
X\arrow[r,"{e_q^X}","{\substack{\text{hom uni}\\\text{de degré }q}}"'] \& \Gamma^q(X)\arrow[r,"{e_p^{\Gamma^q(X)}}","{\substack{\text{hom uni}\\\text{de degré }p.}}"'] \& \Gamma^p(\Gamma^q(X))
\end{tikzcd}}

On a la condition ``d'associativité''
\diagram{$\unsure{(*)|(\times)}\quad$\begin{tikzcd}[ampersand replacement=\&,column sep=large]
{\Gamma^{pqr}(X)}\arrow[r,dash,"{\alpha^X_{p,qr}}"]\arrow[d,"{\alpha^X_{pq,r}}"'] \& {\Gamma^p(\Gamma^{qr}(X))}\arrow[d,"{\Gamma^p(\alpha_{q,r}^{\unsure{\Gamma^r(X)|X}})}"]\\
{\Gamma^{pq}(\Gamma^r(X))}\arrow[r,"{\deleted{\unsure{\Gamma}}\alpha^{\Gamma^r(X)}_{p,q}}"] \& {\Gamma^p(\Gamma^q(\Gamma^r(X)))}
\arrow[from=1-1,to=2-2,phantom,"\text{comm.}" description]
\end{tikzcd}}
et d'\uline{unité} : l'isom \unsure{$\alpha$|.} est l'identité \unsure{si}
\unsure{l'on pose|l'on prend} $\Gamma^1 X=X$, $\Gamma^0 X=\mathbb{1}$
% archive sheet number: 80
